\begin{solapartat}
\[ \vec{F} = q\,\vec{v}\wedge\vec{B} \ \text{\punts{0,3}} = 1,6\cdot 10^{-19}\times 2\cdot 10^{6}\times 0,5\,(\vec{\imath}\wedge\vec{\jmath}) \ \text{\punts{0,3}} = 1,6\cdot 10^{-13}\,\vec{k}\ \mathrm{N} \ \text{\punts{0,2}} \]
Per tant la força va dirigida en sentit vertical \punts{0,2}

En el cas en que no donin correctament la direcció de la força, restarem \textbf{0,2} punts.
\end{solapartat}

\begin{solapartat}
Al ser la força perpendicular a la velocitat, el moviment serà el d'un moviment circular uniforme \punts{0,2}

La força que fa el camp magnètic sobre els protons es la que proporciona l'acceleració centrípeta que farà girar els protons, per trobar el radi de la trajectòria circular tindrem:
\[ q\,v\,B = m\,\frac{v^2}{r} \ \text{\punts{0,2}} \Rightarrow r = \frac{m\,v}{q\,B} \ \text{\punts{0,3}} = 4,18\cdot 10^{-2}\ \mathrm{m} \ \text{\punts{0,2}} \]
Els protons no impactaran ningú, ja que només d'entrar l'aula fan la trajectòria circular de radi $4,18\cdot 10^{-2}\ \mathrm{m}$ \punts{0,1}
\end{solapartat}
